\documentclass[aspectratio=169,11pt]{beamer}

\usepackage[utf8]{inputenc}
\usepackage[T1]{fontenc}
\usepackage[spanish,es-tabla]{babel}
\usepackage{booktabs}
\usepackage{tabularx}
\usepackage{array}
\usepackage{amsmath}
\usepackage{xcolor}
\usepackage{hyperref}

\usetheme{Madrid}
\usecolortheme{default}
\definecolor{arrozverde}{RGB}{34,105,82}
\definecolor{arrozoscuro}{RGB}{31,49,61}
\definecolor{arrozoro}{RGB}{196,143,45}
\setbeamercolor{structure}{fg=arrozverde}
\setbeamercolor{title}{fg=arrozoscuro}
\setbeamercolor{frametitle}{fg=arrozoscuro}
\setbeamercolor{alerted text}{fg=arrozoro}
\setbeamertemplate{navigation symbols}{}
\setbeamertemplate{footline}[frame number]
\newcommand{\tiempo}[1]{\vfill\hfill\scriptsize\textit{Tiempo sugerido: #1}}
\newcommand{\pendiente}{\textcolor{arrozoro}{\textbf{pendiente}}}

\title{Comparación de modelos para clasificar enfermedades en hojas de arroz}
\subtitle{Paddy Doctor: CNN, Vision Transformer e híbrido}
\author{[Nombre del autor]}
\institute{[Institución]}
\date{\today}

\begin{document}

\begin{frame}
    \titlepage
    \begin{center}
        \small Exposición estimada: 10 minutos \textbullet\ Síntesis bibliográfica y benchmark local en curso
    \end{center}
    \tiempo{0:20}
\end{frame}

\begin{frame}{Problema y pregunta}
    \begin{itemize}
        \item Las enfermedades del arroz afectan rendimiento y calidad; una foto puede apoyar una evaluación rápida.
        \item \textbf{Pregunta:} ¿qué arquitectura reconoce mejor las diez clases del conjunto local sin ignorar las clases minoritarias?
        \item No se asume que el modelo más complejo sea el mejor: mediremos precisión y costo.
        \item El objetivo es \textbf{clasificar una imagen completa}; no localizar lesiones ni sustituir al especialista.
    \end{itemize}
    \tiempo{0:50}
\end{frame}

\begin{frame}{Qué aportan los estudios previos}
    \small
    \begin{tabularx}{\textwidth}{p{3.3cm}p{4.0cm}Y}
        \toprule
        Estudio & Propuesta y cifra reportada & Cómo interpretarla aquí \\
        \midrule
        Petchiammal et al.\ (2023) & ResNet34; F1 97.50\%; 16\,225 imágenes y 13 clases. & Referencia sobre Paddy Doctor; no es la misma partición local. \\
        Khanal et al.\ (2024) & ViT: accuracy 99.38\%; YOLOv8: mAP@0.5 69\%. & Detección y clasificación son tareas diferentes; el artículo tiene cifras internas discrepantes. \\
        Subbarayudu y Kubendiran (2025) & 13\,866 imágenes/13 clases; MobileNetV3 + GKSO-SA + CatBoost; accuracy 98.52\%. & Pipeline diferente; el artículo reporta también F1 97.54\%. \\
        \bottomrule
    \end{tabularx}
    \vspace{0.25cm}
    \begin{block}{Idea clave}
        Son antecedentes, no una tabla de posiciones directa: cambian el número de clases, la versión de datos y el protocolo.
    \end{block}
    \tiempo{1:05}
\end{frame}

\begin{frame}{Nuestro conjunto de datos}
    \begin{columns}[T]
        \column{0.48\textwidth}
        \begin{block}{Tamaño y partición}
            \begin{itemize}
                \item 10\,407 imágenes, 10 clases.
                \item Train: 8\,325 (80\%).
                \item Validación: 1\,041 (10\%).
                \item Test: 1\,041 (10\%).
                \item División estratificada, semilla 42, idéntica para los modelos.
            \end{itemize}
        \end{block}
        \column{0.48\textwidth}
        \begin{block}{Desbalance}
            \begin{itemize}
                \item Ejemplos: bacterial panicle blight 337; bacterial leaf streak 380.
                \item Clases más frecuentes: normal 1\,764; blast 1\,738.
                \item Se pondera la pérdida usando solo las frecuencias de train.
            \end{itemize}
        \end{block}
    \end{columns}
    \tiempo{0:55}
\end{frame}

\begin{frame}{Metodología: datos y aumentos}
    \begin{enumerate}
        \item Imágenes RGB a \(224\times224\); normalización ImageNet.
        \item Aumentos \textbf{solo en train}, aleatorios en cada época:
        \begin{itemize}
            \item volteos horizontal/vertical;
            \item pequeñas rotaciones, traslaciones, cambios de escala y shear;
            \item variaciones moderadas de brillo, contraste, saturación y tono.
        \end{itemize}
        \item Validación/test: solo redimensionamiento y normalización.
        \item El aumento es en línea: crea variantes temporales, no nuevas copias en disco.
    \end{enumerate}
    \begin{alertblock}{Prevención de fuga}
        La partición se realiza antes del aumento; el test se reserva y se consulta al final.
    \end{alertblock}
    \tiempo{1:20}
\end{frame}

\begin{frame}{Cuatro modelos, un protocolo}
    \begin{columns}[T]
        \column{0.49\textwidth}
        \begin{block}{Modelos individuales}
            \begin{itemize}
                \item \textbf{ResNet34}: CNN residual; referencia del benchmark.
                \item \textbf{EfficientNet-B0}: CNN eficiente; baseline compacto.
                \item \textbf{ViT-Small/16}: Transformer visual; atención global.
            \end{itemize}
        \end{block}
        \column{0.49\textwidth}
        \begin{block}{Híbrido}
            \textbf{EfficientNet-B0 + ViT-Small/16}
            \begin{itemize}
                \item Extrae rasgos con las dos ramas preentrenadas.
                \item Concatena rasgos globales.
                \item Cabeza común: LayerNorm, dropout y clasificador.
            \end{itemize}
        \end{block}
    \end{columns}
    \vspace{0.2cm}
    Todos parten de pesos ImageNet y usan las mismas particiones e imágenes.
    \tiempo{1:05}
\end{frame}

\begin{frame}{Entrenamiento y criterios de evaluación}
    \begin{itemize}
        \item Transfer learning con ajuste fino; AdamW, tasas diferenciadas para backbone y cabeza.
        \item Batch 16; máximo 20 épocas; early stopping tras 4 épocas sin mejora.
        \item Checkpoint elegido por \textbf{F1 macro de validación}; precisión mixta solo si CUDA está disponible.
        \item Métrica principal: \textbf{F1 macro}; complementarias: accuracy, balanced accuracy y F1 ponderado.
        \item Reporte por clase y matriz de confusión para ver qué enfermedades se confunden.
    \end{itemize}
    \begin{block}{Regla de decisión}
        Comparar primero F1 macro y recall por clase; luego considerar memoria, tiempo e inferencia.
    \end{block}
    \tiempo{1:10}
\end{frame}

\begin{frame}{Referencias provisionales: resultados publicados}
    \small
    \begin{tabularx}{\textwidth}{p{3.1cm}p{4.0cm}Y}
            \toprule
            Estudio & Modelo o pipeline & Cifras reportadas por los autores \\
            \midrule
            Petchiammal et al.\ (2023) & ResNet34; 16\,225 imágenes/13 clases & F1: 97.50\% \\
            Khanal et al.\ (2024) & Vision Transformer; el mismo artículo también evalúa YOLOv8 & Accuracy: 99.38\%; YOLOv8 mAP@0.5: 69\% \\
            Subbarayudu y Kubendiran (2025) & MobileNetV3 + GKSO-SA + CatBoost; 13\,866 imágenes/13 clases & Accuracy: 98.52\%; F1: 97.54\% \\
            \bottomrule
    \end{tabularx}
    \begin{alertblock}{Atribución}
        Estas son cifras publicadas por los tres estudios; no son resultados de nuestros entrenamientos.
    \end{alertblock}
    \tiempo{1:10}
\end{frame}

\begin{frame}{Avance actual y lectura comparativa}
    \begin{itemize}
        \item \textbf{Captura de la corrida}: ResNet34 terminó 20/20; F1 macro val. 0.9418 y accuracy val. 0.9452 en la última época visible.
        \item EfficientNet-B0 va por 3/20: F1 macro val. 0.8567 y accuracy val. 0.8809; su métrica aún puede cambiar.
        \item La captura no muestra dispositivo y no incluye test; son avances, no resultados finales.
        \item En los artículos, la mayor accuracy publicada es la del Vision Transformer: 99.38\%. Una posible ventaja es el contexto global, pero no es prueba de causalidad.
        \item \textbf{No hay ganador demostrado entre estudios}: cambian datos y protocolos; la brecha 99.38--98.52 es descriptiva, no una comparación controlada.
    \end{itemize}
    \begin{block}{Resultado propio}
        Faltan ViT-Small/16, el híbrido y las métricas de test para elegir el mejor en nuestro conjunto.
    \end{block}
    \tiempo{0:55}
\end{frame}

\begin{frame}{Conclusiones y próximos pasos}
    \begin{enumerate}
        \item Terminar el benchmark local y registrar dispositivo, versiones y checkpoints.
        \item Actualizar las métricas de test y revisar recall/F1 por clase y matrices de confusión.
        \item Elegir el modelo por equilibrio entre desempeño y costo; no extrapolar el 99.38\% de otro estudio.
        \item Validar después con fotos externas y, si el uso lo requiere, añadir explicabilidad y evaluación agronómica.
    \end{enumerate}
    \begin{block}{Conclusión empírica}
        En la literatura, ViT tiene la mayor accuracy reportada; en nuestro conjunto, el ganador sigue \pendiente.
    \end{block}
    \tiempo{1:00}
\end{frame}

\begin{frame}{Referencias breves}
    \footnotesize
    \begin{itemize}
        \item Petchiammal et al.\ (2023). \emph{Paddy Doctor: A Visual Image Dataset for Automated Paddy Disease Classification and Benchmarking}. DOI: \url{https://doi.org/10.1145/3570991.3570994}.
        \item Khanal et al.\ (2024). \emph{Paddy Disease Detection and Classification Using Computer Vision Techniques}. arXiv:2412.05996.
        \item Subbarayudu y Kubendiran (2025). \emph{An automated hybrid deep learning framework for paddy leaf disease identification and classification}. \emph{Scientific Reports}, 15. DOI: \url{https://doi.org/10.1038/s41598-025-08071-6}.
    \end{itemize}
    \begin{center}
        \large Gracias
    \end{center}
    \tiempo{0:10}
\end{frame}

\end{document}
